\section{高宗时期的工程、知识与区域生态}
高宗时期的河工、矿业、盐政、屯田、驻防、书籍征集与高原交通，使制度史与环境史无法分开。本节不以国家工程或文教项目作为成功的同义词，而以劳役、材料、季节、运输、翻译和地方知识的可见与不可见为中心。
\subsection{1720—1729：区域网络、制度媒介与环境条件}
这些事件的先后可以由账本定位，但每一项影响都须回到具体区域、材料与行动者，不能被压缩为抽象的国家趋势。
\hypertarget{civic:EQ-1723-YONGZHENG-ACCESSION}{}\paragraph{雍正元年（1723年）}雍正帝即位并开始重整康熙晚年的财政、宗室与官僚议题。本事件的制度含义须同工程、运输、劳役、生态条件或知识传播的实际媒介结合，不能把一道制度或一部汇编当成地方实践本身。 核心取证为《清世宗实录》雍正元年相关条；宫中朱批奏折、内阁大库题本与《清史稿》相关志传。行政文书显示的是机构如何把道路、港口、人口、宗教或案件转换为可处理的类别；没有后续县档或物质证据，不能由分类推定实际生活。账本提示的研究线索为宫廷政治、制度运作与中央—地方信息链研究。事件所处条件为康熙帝去世后继承程序与皇子网络需要迅速稳定。后续变化应限于新帝以密折、用人与财政整顿重塑决策中心。来源限制：以雍正元年（1723年）的同期文书为定位起点；官修记录、奏报或外交文本服务特定机构立场，具体日期、规模、地方执行与对手视角须以独立材料复核。
\hypertarget{civic:EQ-1723-JI-ZENGYUN-CASE}{}\paragraph{雍正元年（1723年）}年羹尧在西北军政中的权力与朝廷猜忌加剧。本事件的制度含义须同工程、运输、劳役、生态条件或知识传播的实际媒介结合，不能把一道制度或一部汇编当成地方实践本身。 核心取证为《清世宗实录》雍正元年相关条；宫中朱批奏折、内阁大库题本与《清史稿》相关志传。行政文书显示的是机构如何把道路、港口、人口、宗教或案件转换为可处理的类别；没有后续县档或物质证据，不能由分类推定实际生活。账本提示的研究线索为宫廷政治、制度运作与中央—地方信息链研究。事件所处条件为边军统帅兼具军功、财政与人事资源。后续变化应限于其后对年氏集团的处分成为雍正整顿权力的节点。来源限制：以雍正元年（1723年）的同期文书为定位起点；官修记录、奏报或外交文本服务特定机构立场，具体日期、规模、地方执行与对手视角须以独立材料复核。
\hypertarget{civic:EQ-1724-CATHOLIC-PROSCRIPTION}{}\paragraph{雍正二年（1724年）}清廷扩大对天主教活动的限制，将传教、地方教民与外来人员置入治理问题。本事件的制度含义须同工程、运输、劳役、生态条件或知识传播的实际媒介结合，不能把一道制度或一部汇编当成地方实践本身。 核心取证为《清世宗实录》雍正二年相关条；地方志、督抚奏折及地方档案。编年条目可以为行动建立相对次序，却通常不保存劳动过程、物价变动、建筑材料或非精英的叙述。账本提示的研究线索为地方档案、迁徙、宗教、族群与社会动员研究。事件所处条件为礼仪争论与地方官对宗教网络的疑虑相互叠加。后续变化应限于宗教实践转入不同地区和身份条件下的隐蔽或协商状态。来源限制：以雍正二年（1724年）的同期文书为定位起点；官修记录、奏报或外交文本服务特定机构立场，具体日期、规模、地方执行与对手视角须以独立材料复核。
\hypertarget{civic:EQ-1725-NIAN-GENGYAO-PUNISHMENT}{}\paragraph{雍正三年（1725年）}年羹尧被削职、赐死，其党羽与相关官员受到清理。本事件的制度含义须同工程、运输、劳役、生态条件或知识传播的实际媒介结合，不能把一道制度或一部汇编当成地方实践本身。 核心取证为《清世宗实录》雍正三年相关条；宫中朱批奏折、内阁大库题本与《清史稿》相关志传。行政文书显示的是机构如何把道路、港口、人口、宗教或案件转换为可处理的类别；没有后续县档或物质证据，不能由分类推定实际生活。账本提示的研究线索为宫廷政治、制度运作与中央—地方信息链研究。事件所处条件为雍正帝须限制功臣和边帅形成独立政治中心。后续变化应限于中央对军政人事和密折信息链的控制增强。来源限制：以雍正三年（1725年）的同期文书为定位起点；官修记录、奏报或外交文本服务特定机构立场，具体日期、规模、地方执行与对手视角须以独立材料复核。
\hypertarget{civic:EQ-1725-TAN-TEBI-CASE}{}\paragraph{雍正三年（1725年）}隆科多失势并被幽禁，雍正初年近臣权力结构继续调整。本事件的制度含义须同工程、运输、劳役、生态条件或知识传播的实际媒介结合，不能把一道制度或一部汇编当成地方实践本身。 核心取证为《清世宗实录》雍正三年相关条；宫中朱批奏折、内阁大库题本与《清史稿》相关志传。行政文书显示的是机构如何把道路、港口、人口、宗教或案件转换为可处理的类别；没有后续县档或物质证据，不能由分类推定实际生活。账本提示的研究线索为宫廷政治、制度运作与中央—地方信息链研究。事件所处条件为继承过程中的近臣角色引发君主对安全与忠诚的担忧。后续变化应限于内廷、宗室与部院之间的政治边界被重新划定。来源限制：以雍正三年（1725年）的同期文书为定位起点；官修记录、奏报或外交文本服务特定机构立场，具体日期、规模、地方执行与对手视角须以独立材料复核。
\hypertarget{civic:EQ-1726-HAOXIAN-GUIGONG}{}\paragraph{雍正四年（1726年）}雍正政府推进耗羡归公与养廉银安排，试图重组地方财政和官员俸给。本事件的制度含义须同工程、运输、劳役、生态条件或知识传播的实际媒介结合，不能把一道制度或一部汇编当成地方实践本身。 核心取证为《清世宗实录》雍正四年相关条；户部奏销、盐政、河工或海关档案。行政文书显示的是机构如何把道路、港口、人口、宗教或案件转换为可处理的类别；没有后续县档或物质证据，不能由分类推定实际生活。账本提示的研究线索为财政账册、市场网络、灾害环境与区域经济研究。事件所处条件为正额俸银不足与附加征收长期并存。后续变化应限于地方财政的公开名目扩大，但实际征收和地区差异仍待账册核验。来源限制：以雍正四年（1726年）的同期文书为定位起点；官修记录、奏报或外交文本服务特定机构立场，具体日期、规模、地方执行与对手视角须以独立材料复核。
\hypertarget{civic:EQ-1726-Gaitu-Guiliu-SOUTHWEST}{}\paragraph{雍正四年（1726年）}清廷在西南继续改土归流，将部分土司辖地改置流官体系。本事件的制度含义须同工程、运输、劳役、生态条件或知识传播的实际媒介结合，不能把一道制度或一部汇编当成地方实践本身。 核心取证为《清世宗实录》雍正四年相关条；军机处档、理藩院档或相关方略。编年条目可以为行动建立相对次序，却通常不保存劳动过程、物价变动、建筑材料或非精英的叙述。账本提示的研究线索为多语边疆档案、盟旗/寺院/绿洲地方史与帝国治理研究。事件所处条件为矿产、交通、纠纷与军事控制推动中央介入。后续变化应限于行政改置伴随地方抵抗、迁徙和长期治理成本。来源限制：以雍正四年（1726年）的同期文书为定位起点；官修记录、奏报或外交文本服务特定机构立场，具体日期、规模、地方执行与对手视角须以独立材料复核。
\hypertarget{civic:EQ-1727-KYAKHTA-TREATY}{}\paragraph{雍正五年（1727年）}《恰克图条约》确定清俄北部边界与贸易、使团安排。本事件的制度含义须同工程、运输、劳役、生态条件或知识传播的实际媒介结合，不能把一道制度或一部汇编当成地方实践本身。 核心取证为《清世宗实录》雍正五年相关条；《筹办夷务始末》相关年卷与中外照会、条约文本。行政文书显示的是机构如何把道路、港口、人口、宗教或案件转换为可处理的类别；没有后续县档或物质证据，不能由分类推定实际生活。账本提示的研究线索为条约文本、外交档案、口岸社会与国际法实践研究。事件所处条件为边界冲突和蒙古地区秩序需要制度化协商。后续变化应限于恰克图成为跨境贸易与外交往来的重要节点。来源限制：以雍正五年（1727年）的同期文书为定位起点；官修记录、奏报或外交文本服务特定机构立场，具体日期、规模、地方执行与对手视角须以独立材料复核。
\hypertarget{civic:EQ-1727-BANNER-GENEALOGIES}{}\paragraph{雍正五年（1727年）}清廷命令编修八旗谱牒，试图界定旗籍、世系与待遇资格。本事件的制度含义须同工程、运输、劳役、生态条件或知识传播的实际媒介结合，不能把一道制度或一部汇编当成地方实践本身。 核心取证为《清世宗实录》雍正五年相关条；地方志、督抚奏折及地方档案。编年条目可以为行动建立相对次序，却通常不保存劳动过程、物价变动、建筑材料或非精英的叙述。账本提示的研究线索为地方档案、迁徙、宗教、族群与社会动员研究。事件所处条件为旗籍资源、婚姻与身份管理需要可查的分类文书。后续变化应限于旗人身份被更紧密地纳入国家档案和待遇秩序。来源限制：以雍正五年（1727年）的同期文书为定位起点；官修记录、奏报或外交文本服务特定机构立场，具体日期、规模、地方执行与对手视角须以独立材料复核。
\hypertarget{civic:EQ-1728-Dzungar-EXPEDITION}{}\paragraph{雍正六年（1728年）}清军在西北对准噶尔展开远征，后勤与高原草原路线成为军事难题。本事件的制度含义须同工程、运输、劳役、生态条件或知识传播的实际媒介结合，不能把一道制度或一部汇编当成地方实践本身。 核心取证为《清世宗实录》雍正六年相关条；军机处档、战事方略及地方奏报。编年条目可以为行动建立相对次序，却通常不保存劳动过程、物价变动、建筑材料或非精英的叙述。账本提示的研究线索为战争动员、军需、战场暴力与地方社会研究。事件所处条件为清廷希望限制准噶尔对蒙古与青海的影响。后续变化应限于军需困难促使中央反复调整将领、补给和联盟策略。来源限制：以雍正六年（1728年）的同期文书为定位起点；官修记录、奏报或外交文本服务特定机构立场，具体日期、规模、地方执行与对手视角须以独立材料复核。
\hypertarget{civic:EQ-1728-ZHUNGHAR-TIBET-SETTLEMENT}{}\paragraph{雍正六年（1728年）}清廷在准噶尔威胁下调整驻藏与青海蒙古的军事部署。本事件的制度含义须同工程、运输、劳役、生态条件或知识传播的实际媒介结合，不能把一道制度或一部汇编当成地方实践本身。 核心取证为《清世宗实录》雍正六年相关条；军机处档、理藩院档或相关方略。编年条目可以为行动建立相对次序，却通常不保存劳动过程、物价变动、建筑材料或非精英的叙述。账本提示的研究线索为多语边疆档案、盟旗/寺院/绿洲地方史与帝国治理研究。事件所处条件为拉萨政局和草原力量关系持续不稳。后续变化应限于西藏事务被更直接地纳入清廷边务协调。来源限制：以雍正六年（1728年）的同期文书为定位起点；官修记录、奏报或外交文本服务特定机构立场，具体日期、规模、地方执行与对手视角须以独立材料复核。
\hypertarget{civic:EQ-1729-JUNJICHU-FOUNDING}{}\paragraph{雍正七年（1729年）}雍正帝设军机房，后来发展为军机处，处理紧急军政文书。本事件的制度含义须同工程、运输、劳役、生态条件或知识传播的实际媒介结合，不能把一道制度或一部汇编当成地方实践本身。 核心取证为《清世宗实录》雍正七年相关条；宫中朱批奏折、内阁大库题本与《清史稿》相关志传。行政文书显示的是机构如何把道路、港口、人口、宗教或案件转换为可处理的类别；没有后续县档或物质证据，不能由分类推定实际生活。账本提示的研究线索为宫廷政治、制度运作与中央—地方信息链研究。事件所处条件为西北战争要求比常规部院更快的保密决策链。后续变化应限于军机处成为中后期中央决策的重要枢纽。来源限制：以雍正七年（1729年）的同期文书为定位起点；官修记录、奏报或外交文本服务特定机构立场，具体日期、规模、地方执行与对手视角须以独立材料复核。
\subsection{1730—1739：区域网络、制度媒介与环境条件}
这些事件的先后可以由账本定位，但每一项影响都须回到具体区域、材料与行动者，不能被压缩为抽象的国家趋势。
\hypertarget{civic:EQ-1730-SOUTHWEST-MIAO-GOVERNANCE}{}\paragraph{雍正八年（1730年）}清廷围绕苗疆设置、屯兵与地方纠纷加强治理。本事件的制度含义须同工程、运输、劳役、生态条件或知识传播的实际媒介结合，不能把一道制度或一部汇编当成地方实践本身。 核心取证为《清世宗实录》雍正八年相关条；军机处档、理藩院档或相关方略。编年条目可以为行动建立相对次序，却通常不保存劳动过程、物价变动、建筑材料或非精英的叙述。账本提示的研究线索为多语边疆档案、盟旗/寺院/绿洲地方史与帝国治理研究。事件所处条件为改土归流与移民开发改变既有地方权力关系。后续变化应限于军政介入与社区抵抗交织，不能以行政设置等同稳定统治。来源限制：以雍正八年（1730年）的同期文书为定位起点；官修记录、奏报或外交文本服务特定机构立场，具体日期、规模、地方执行与对手视角须以独立材料复核。
\hypertarget{civic:EQ-1731-Dzungar-WAR-REVERSE}{}\paragraph{雍正九年（1731年）}清军在对准噶尔战争中受挫，西北远征暴露补给与指挥困难。本事件的制度含义须同工程、运输、劳役、生态条件或知识传播的实际媒介结合，不能把一道制度或一部汇编当成地方实践本身。 核心取证为《清世宗实录》雍正九年相关条；军机处档、战事方略及地方奏报。编年条目可以为行动建立相对次序，却通常不保存劳动过程、物价变动、建筑材料或非精英的叙述。账本提示的研究线索为战争动员、军需、战场暴力与地方社会研究。事件所处条件为长距离运输、气候和草原盟友关系限制军队行动。后续变化应限于朝廷调整统帅、粮道与军费安排。来源限制：以雍正九年（1731年）的同期文书为定位起点；官修记录、奏报或外交文本服务特定机构立场，具体日期、规模、地方执行与对手视角须以独立材料复核。
\hypertarget{civic:EQ-1732-KOBDO-BATTLE}{}\paragraph{雍正十年（1732年）}清军在和通泊等战事失利，西北战局一度紧张。本事件的制度含义须同工程、运输、劳役、生态条件或知识传播的实际媒介结合，不能把一道制度或一部汇编当成地方实践本身。 核心取证为《清世宗实录》雍正十年相关条；军机处档、战事方略及地方奏报。编年条目可以为行动建立相对次序，却通常不保存劳动过程、物价变动、建筑材料或非精英的叙述。账本提示的研究线索为战争动员、军需、战场暴力与地方社会研究。事件所处条件为准噶尔骑兵和草原地理使清军难以速胜。后续变化应限于清廷转向重整蒙古盟友、驻防与运输体系。来源限制：以雍正十年（1732年）的同期文书为定位起点；官修记录、奏报或外交文本服务特定机构立场，具体日期、规模、地方执行与对手视角须以独立材料复核。
\hypertarget{civic:EQ-1732-AMIN-KHOJA-FLIGHT}{}\paragraph{雍正十年（1732年）}阿敏和卓等人在准噶尔压力下向清境迁移，绿洲政治进入清准竞争。本事件的制度含义须同工程、运输、劳役、生态条件或知识传播的实际媒介结合，不能把一道制度或一部汇编当成地方实践本身。 核心取证为《清世宗实录》雍正十年相关条；军机处档、理藩院档或相关方略。编年条目可以为行动建立相对次序，却通常不保存劳动过程、物价变动、建筑材料或非精英的叙述。账本提示的研究线索为多语边疆档案、盟旗/寺院/绿洲地方史与帝国治理研究。事件所处条件为准噶尔内部统治与绿洲地方关系紧张。后续变化应限于清廷获得介入天山南路事务的新信息和人际网络。来源限制：以雍正十年（1732年）的同期文书为定位起点；官修记录、奏报或外交文本服务特定机构立场，具体日期、规模、地方执行与对手视角须以独立材料复核。
\hypertarget{civic:EQ-1733-YONGZHENG-FISCAL-INSPECTION}{}\paragraph{雍正十一年（1733年）}雍正政府持续以奏销、亏空追查和官员考成为手段整顿地方财政。本事件的制度含义须同工程、运输、劳役、生态条件或知识传播的实际媒介结合，不能把一道制度或一部汇编当成地方实践本身。 核心取证为《清世宗实录》雍正十一年相关条；户部奏销、盐政、河工或海关档案。行政文书显示的是机构如何把道路、港口、人口、宗教或案件转换为可处理的类别；没有后续县档或物质证据，不能由分类推定实际生活。账本提示的研究线索为财政账册、市场网络、灾害环境与区域经济研究。事件所处条件为耗羡改革需要与旧有钱粮、仓储和欠额制度衔接。后续变化应限于地方官的呈报策略与财政实际之间形成新的张力。来源限制：以雍正十一年（1733年）的同期文书为定位起点；官修记录、奏报或外交文本服务特定机构立场，具体日期、规模、地方执行与对手视角须以独立材料复核。
\hypertarget{civic:EQ-1734-CHRISTIAN-ENFORCEMENT}{}\paragraph{雍正十二年（1734年）}地方官继续执行限制天主教的政策，传教士和教民的活动受地域差异影响。本事件的制度含义须同工程、运输、劳役、生态条件或知识传播的实际媒介结合，不能把一道制度或一部汇编当成地方实践本身。 核心取证为《清世宗实录》雍正十二年相关条；地方志、督抚奏折及地方档案。编年条目可以为行动建立相对次序，却通常不保存劳动过程、物价变动、建筑材料或非精英的叙述。账本提示的研究线索为地方档案、迁徙、宗教、族群与社会动员研究。事件所处条件为中央禁令需经省县转化为具体案件与登记。后续变化应限于不同地区的执行强度不能由中央谕旨一概而论。来源限制：以雍正十二年（1734年）的同期文书为定位起点；官修记录、奏报或外交文本服务特定机构立场，具体日期、规模、地方执行与对手视角须以独立材料复核。
\hypertarget{civic:EQ-1735-QIANLONG-ACCESSION}{}\paragraph{雍正十三年（1735年）}乾隆帝即位，雍正时期形成的财政与军机决策机制进入新朝。本事件的制度含义须同工程、运输、劳役、生态条件或知识传播的实际媒介结合，不能把一道制度或一部汇编当成地方实践本身。 核心取证为《清世宗实录》雍正十三年相关条；宫中朱批奏折、内阁大库题本与《清史稿》相关志传。行政文书显示的是机构如何把道路、港口、人口、宗教或案件转换为可处理的类别；没有后续县档或物质证据，不能由分类推定实际生活。账本提示的研究线索为宫廷政治、制度运作与中央—地方信息链研究。事件所处条件为雍正帝去世与皇位继承完成。后续变化应限于乾隆初政开始重估前朝用人、刑狱与边务方针。来源限制：以雍正十三年（1735年）的同期文书为定位起点；官修记录、奏报或外交文本服务特定机构立场，具体日期、规模、地方执行与对手视角须以独立材料复核。
\hypertarget{civic:EQ-1735-MIAO-UPRISING-RONGJIANG}{}\paragraph{雍正十三年（1735年）}贵州荣江等地苗侗社区反抗清军和地方治理，镇压造成严重暴力。本事件的制度含义须同工程、运输、劳役、生态条件或知识传播的实际媒介结合，不能把一道制度或一部汇编当成地方实践本身。 核心取证为《清世宗实录》雍正十三年相关条；军机处档、战事方略及地方奏报。编年条目可以为行动建立相对次序，却通常不保存劳动过程、物价变动、建筑材料或非精英的叙述。账本提示的研究线索为战争动员、军需、战场暴力与地方社会研究。事件所处条件为行政扩张、赋役与地方资源冲突累积。后续变化应限于官方战报的伤亡数字和族群标签须与地方材料互校。来源限制：以雍正十三年（1735年）的同期文书为定位起点；官修记录、奏报或外交文本服务特定机构立场，具体日期、规模、地方执行与对手视角须以独立材料复核。
\hypertarget{civic:EQ-1736-AMNESTY-AND-RECTIFICATION}{}\paragraph{乾隆元年（1736年）}乾隆初年颁布恩诏并调整雍正后期政治案件的处置。本事件的制度含义须同工程、运输、劳役、生态条件或知识传播的实际媒介结合，不能把一道制度或一部汇编当成地方实践本身。 核心取证为《清高宗实录》乾隆元年相关条；宫中朱批奏折、内阁大库题本与《清史稿》相关志传。行政文书显示的是机构如何把道路、港口、人口、宗教或案件转换为可处理的类别；没有后续县档或物质证据，不能由分类推定实际生活。账本提示的研究线索为宫廷政治、制度运作与中央—地方信息链研究。事件所处条件为新君需要建立自身施政形象并重组官僚关系。后续变化应限于部分案件被重新解释，但中央集权并未放松。来源限制：以乾隆元年（1736年）的同期文书为定位起点；官修记录、奏报或外交文本服务特定机构立场，具体日期、规模、地方执行与对手视角须以独立材料复核。
\hypertarget{civic:EQ-1737-Dzungar-NEGOTIATION}{}\paragraph{乾隆二年（1737年）}清廷在战争压力下继续处理与准噶尔的使节、盟旗和军事情报。本事件的制度含义须同工程、运输、劳役、生态条件或知识传播的实际媒介结合，不能把一道制度或一部汇编当成地方实践本身。 核心取证为《清高宗实录》乾隆二年相关条；《筹办夷务始末》相关年卷与中外照会、条约文本。行政文书显示的是机构如何把道路、港口、人口、宗教或案件转换为可处理的类别；没有后续县档或物质证据，不能由分类推定实际生活。账本提示的研究线索为条约文本、外交档案、口岸社会与国际法实践研究。事件所处条件为草原战争难以只靠单线军事行动解决。后续变化应限于外交、情报与军事准备并行发展。来源限制：以乾隆二年（1737年）的同期文书为定位起点；官修记录、奏报或外交文本服务特定机构立场，具体日期、规模、地方执行与对手视角须以独立材料复核。
\hypertarget{civic:EQ-1738-FIRST-JINCHUAN-WAR}{}\paragraph{乾隆三年（1738年）}清军开始大规模干预大金川，山地战与土司关系限制军事推进。本事件的制度含义须同工程、运输、劳役、生态条件或知识传播的实际媒介结合，不能把一道制度或一部汇编当成地方实践本身。 核心取证为《清高宗实录》乾隆三年相关条；军机处档、战事方略及地方奏报。编年条目可以为行动建立相对次序，却通常不保存劳动过程、物价变动、建筑材料或非精英的叙述。账本提示的研究线索为战争动员、军需、战场暴力与地方社会研究。事件所处条件为地方权力冲突与清廷对川西治理的扩展相互作用。后续变化应限于远征显示高山补给和地方联盟的成本极高。来源限制：以乾隆三年（1738年）的同期文书为定位起点；官修记录、奏报或外交文本服务特定机构立场，具体日期、规模、地方执行与对手视角须以独立材料复核。
\hypertarget{civic:EQ-1739-JINCHUAN-SUPPLY}{}\paragraph{乾隆四年（1739年）}大金川战事的粮运、道路和军费成为中央持续关注的难题。本事件的制度含义须同工程、运输、劳役、生态条件或知识传播的实际媒介结合，不能把一道制度或一部汇编当成地方实践本身。 核心取证为《清高宗实录》乾隆四年相关条；户部奏销、盐政、河工或海关档案。行政文书显示的是机构如何把道路、港口、人口、宗教或案件转换为可处理的类别；没有后续县档或物质证据，不能由分类推定实际生活。账本提示的研究线索为财政账册、市场网络、灾害环境与区域经济研究。事件所处条件为山地驻军需要远距离转运粮械。后续变化应限于军费和民夫负担使战争后果超出战场范围。来源限制：以乾隆四年（1739年）的同期文书为定位起点；官修记录、奏报或外交文本服务特定机构立场，具体日期、规模、地方执行与对手视角须以独立材料复核。
\subsection{1740—1749：区域网络、制度媒介与环境条件}
这些事件的先后可以由账本定位，但每一项影响都须回到具体区域、材料与行动者，不能被压缩为抽象的国家趋势。
\hypertarget{civic:EQ-1740-NORTH-CHINA-FAMINE}{}\paragraph{乾隆五年（1740年）}华北旱灾和粮价问题促使地方奏报与赈济措施增加。本事件的制度含义须同工程、运输、劳役、生态条件或知识传播的实际媒介结合，不能把一道制度或一部汇编当成地方实践本身。 核心取证为《清高宗实录》乾隆五年相关条；地方志、督抚奏折及地方档案。编年条目可以为行动建立相对次序，却通常不保存劳动过程、物价变动、建筑材料或非精英的叙述。账本提示的研究线索为地方档案、迁徙、宗教、族群与社会动员研究。事件所处条件为气候波动、粮食流通和地方仓储共同影响灾情。后续变化应限于灾情、赈济和死亡迁徙应分开记录并核对地域口径。来源限制：以乾隆五年（1740年）的同期文书为定位起点；官修记录、奏报或外交文本服务特定机构立场，具体日期、规模、地方执行与对手视角须以独立材料复核。
\hypertarget{civic:EQ-1741-MIAO-FRONTIER-CAMPAIGN}{}\paragraph{乾隆六年（1741年）}清廷在西南继续以驻兵、设治和军事行动处理苗疆冲突。本事件的制度含义须同工程、运输、劳役、生态条件或知识传播的实际媒介结合，不能把一道制度或一部汇编当成地方实践本身。 核心取证为《清高宗实录》乾隆六年相关条；军机处档、理藩院档或相关方略。编年条目可以为行动建立相对次序，却通常不保存劳动过程、物价变动、建筑材料或非精英的叙述。账本提示的研究线索为多语边疆档案、盟旗/寺院/绿洲地方史与帝国治理研究。事件所处条件为地方行政延伸与社区自治之间矛盾未解。后续变化应限于长期驻防和地方中介成为治理的重要条件。来源限制：以乾隆六年（1741年）的同期文书为定位起点；官修记录、奏报或外交文本服务特定机构立场，具体日期、规模、地方执行与对手视角须以独立材料复核。
\hypertarget{civic:EQ-1742-BANNER-EXIT-ORDER}{}\paragraph{乾隆七年（1742年）}清廷允许部分汉军旗人出旗，调整旗籍与生计安排。本事件的制度含义须同工程、运输、劳役、生态条件或知识传播的实际媒介结合，不能把一道制度或一部汇编当成地方实践本身。 核心取证为《清高宗实录》乾隆七年相关条；地方志、督抚奏折及地方档案。编年条目可以为行动建立相对次序，却通常不保存劳动过程、物价变动、建筑材料或非精英的叙述。账本提示的研究线索为地方档案、迁徙、宗教、族群与社会动员研究。事件所处条件为八旗人口增长和财政供养压力加重。后续变化应限于出旗改变身份和待遇资格，实际执行具有地区差异。来源限制：以乾隆七年（1742年）的同期文书为定位起点；官修记录、奏报或外交文本服务特定机构立场，具体日期、规模、地方执行与对手视角须以独立材料复核。
\hypertarget{civic:EQ-1743-YELLOW-RIVER-WORKS}{}\paragraph{乾隆八年（1743年）}河工官员就黄河堤防、河道和经费持续奏报并实施工程。本事件的制度含义须同工程、运输、劳役、生态条件或知识传播的实际媒介结合，不能把一道制度或一部汇编当成地方实践本身。 核心取证为《清高宗实录》乾隆八年相关条；户部奏销、盐政、河工或海关档案。行政文书显示的是机构如何把道路、港口、人口、宗教或案件转换为可处理的类别；没有后续县档或物质证据，不能由分类推定实际生活。账本提示的研究线索为财政账册、市场网络、灾害环境与区域经济研究。事件所处条件为河患威胁漕运、田赋与沿河社区。后续变化应限于工程预算与实际河势、民工负担不能由定额直接推断。来源限制：以乾隆八年（1743年）的同期文书为定位起点；官修记录、奏报或外交文本服务特定机构立场，具体日期、规模、地方执行与对手视角须以独立材料复核。
\hypertarget{civic:EQ-1744-JESUIT-CALENDAR-SERVICE}{}\paragraph{乾隆九年（1744年）}宫廷继续利用部分西洋传教士的历算与技术服务，同时维持宗教限制。本事件的制度含义须同工程、运输、劳役、生态条件或知识传播的实际媒介结合，不能把一道制度或一部汇编当成地方实践本身。 核心取证为《清高宗实录》乾隆九年相关条；内阁档、文集或报刊相关记录。编年条目可以为行动建立相对次序，却通常不保存劳动过程、物价变动、建筑材料或非精英的叙述。账本提示的研究线索为知识生产、教育、宗教与公共传播研究。事件所处条件为历法、仪器与宫廷知识需求并未与宗教政策完全重合。后续变化应限于技术合作与传教许可须分开考察。来源限制：以乾隆九年（1744年）的同期文书为定位起点；官修记录、奏报或外交文本服务特定机构立场，具体日期、规模、地方执行与对手视角须以独立材料复核。
\hypertarget{civic:EQ-1745-JINCHUAN-NEGOTIATION}{}\paragraph{乾隆十年（1745年）}清廷在大金川战局中尝试通过招抚、盟约和军事压力重组地方关系。本事件的制度含义须同工程、运输、劳役、生态条件或知识传播的实际媒介结合，不能把一道制度或一部汇编当成地方实践本身。 核心取证为《清高宗实录》乾隆十年相关条；军机处档、理藩院档或相关方略。编年条目可以为行动建立相对次序，却通常不保存劳动过程、物价变动、建筑材料或非精英的叙述。账本提示的研究线索为多语边疆档案、盟旗/寺院/绿洲地方史与帝国治理研究。事件所处条件为长期作战难以迅速取得山地控制。后续变化应限于地方首领的降附文书不等于基层社会立即服从。来源限制：以乾隆十年（1745年）的同期文书为定位起点；官修记录、奏报或外交文本服务特定机构立场，具体日期、规模、地方执行与对手视角须以独立材料复核。
\hypertarget{civic:EQ-1746-FIRST-JINCHUAN-SETTLEMENT}{}\paragraph{乾隆十一年（1746年）}第一次金川战争以和议和行政安排告一段落。本事件的制度含义须同工程、运输、劳役、生态条件或知识传播的实际媒介结合，不能把一道制度或一部汇编当成地方实践本身。 核心取证为《清高宗实录》乾隆十一年相关条；军机处档、战事方略及地方奏报。编年条目可以为行动建立相对次序，却通常不保存劳动过程、物价变动、建筑材料或非精英的叙述。账本提示的研究线索为战争动员、军需、战场暴力与地方社会研究。事件所处条件为军费高昂与战场推进困难促使朝廷接受阶段性解决。后续变化应限于后续土司关系并未根除，为再战留下条件。来源限制：以乾隆十一年（1746年）的同期文书为定位起点；官修记录、奏报或外交文本服务特定机构立场，具体日期、规模、地方执行与对手视角须以独立材料复核。
\hypertarget{civic:EQ-1747-SOUTHWEST-MINING-REGULATION}{}\paragraph{乾隆十二年（1747年）}清廷和地方官围绕银铜矿开采、税课与矿工流动加强管理。本事件的制度含义须同工程、运输、劳役、生态条件或知识传播的实际媒介结合，不能把一道制度或一部汇编当成地方实践本身。 核心取证为《清高宗实录》乾隆十二年相关条；户部奏销、盐政、河工或海关档案。行政文书显示的是机构如何把道路、港口、人口、宗教或案件转换为可处理的类别；没有后续县档或物质证据，不能由分类推定实际生活。账本提示的研究线索为财政账册、市场网络、灾害环境与区域经济研究。事件所处条件为矿业吸引流民、资本和武装群体，治安风险上升。后续变化应限于矿务文书可见国家规制，不能代表全部地下经济。来源限制：以乾隆十二年（1747年）的同期文书为定位起点；官修记录、奏报或外交文本服务特定机构立场，具体日期、规模、地方执行与对手视角须以独立材料复核。
\hypertarget{civic:EQ-1748-SONG-AND-CASE}{}\paragraph{乾隆十三年（1748年）}乾隆朝通过大案处分官员，强化对地方文书、钱粮和吏治的监督。本事件的制度含义须同工程、运输、劳役、生态条件或知识传播的实际媒介结合，不能把一道制度或一部汇编当成地方实践本身。 核心取证为《清高宗实录》乾隆十三年相关条；宫中朱批奏折、内阁大库题本与《清史稿》相关志传。行政文书显示的是机构如何把道路、港口、人口、宗教或案件转换为可处理的类别；没有后续县档或物质证据，不能由分类推定实际生活。账本提示的研究线索为宫廷政治、制度运作与中央—地方信息链研究。事件所处条件为中央担忧地方官隐瞒亏空与相互包庇。后续变化应限于案件成为考成与惩戒制度的示范，但不证明全面廉洁。来源限制：以乾隆十三年（1748年）的同期文书为定位起点；官修记录、奏报或外交文本服务特定机构立场，具体日期、规模、地方执行与对手视角须以独立材料复核。
\hypertarget{civic:EQ-1749-JIANGNAN-TAX-REVIEW}{}\paragraph{乾隆十四年（1749年）}江南钱粮、漕运和士绅关系继续在户部与地方奏折中被讨论。本事件的制度含义须同工程、运输、劳役、生态条件或知识传播的实际媒介结合，不能把一道制度或一部汇编当成地方实践本身。 核心取证为《清高宗实录》乾隆十四年相关条；户部奏销、盐政、河工或海关档案。行政文书显示的是机构如何把道路、港口、人口、宗教或案件转换为可处理的类别；没有后续县档或物质证据，不能由分类推定实际生活。账本提示的研究线索为财政账册、市场网络、灾害环境与区域经济研究。事件所处条件为富庶地区的税粮征解和地方负担具有全国财政意义。后续变化应限于中央定额与地方实际征收之间的差距需依账册检验。来源限制：以乾隆十四年（1749年）的同期文书为定位起点；官修记录、奏报或外交文本服务特定机构立场，具体日期、规模、地方执行与对手视角须以独立材料复核。
\subsection{1750—1759：区域网络、制度媒介与环境条件}
这些事件的先后可以由账本定位，但每一项影响都须回到具体区域、材料与行动者，不能被压缩为抽象的国家趋势。
\hypertarget{civic:EQ-1750-LHASA-UPRISING}{}\paragraph{乾隆十五年（1750年）}拉萨发生政治冲突，驻藏大臣和清军介入处置。本事件的制度含义须同工程、运输、劳役、生态条件或知识传播的实际媒介结合，不能把一道制度或一部汇编当成地方实践本身。 核心取证为《清高宗实录》乾隆十五年相关条；军机处档、理藩院档或相关方略。编年条目可以为行动建立相对次序，却通常不保存劳动过程、物价变动、建筑材料或非精英的叙述。账本提示的研究线索为多语边疆档案、盟旗/寺院/绿洲地方史与帝国治理研究。事件所处条件为西藏地方权力、驻防与宗教政治关系不稳。后续变化应限于清廷随后调整驻藏制度，但不能抹去地方行动者的作用。来源限制：以乾隆十五年（1750年）的同期文书为定位起点；官修记录、奏报或外交文本服务特定机构立场，具体日期、规模、地方执行与对手视角须以独立材料复核。
\hypertarget{civic:EQ-1751-TIBET-ORDINANCE}{}\paragraph{乾隆十六年（1751年）}清廷在拉萨事变后重整驻藏大臣和噶厦相关制度安排。本事件的制度含义须同工程、运输、劳役、生态条件或知识传播的实际媒介结合，不能把一道制度或一部汇编当成地方实践本身。 核心取证为《清高宗实录》乾隆十六年相关条；军机处档、理藩院档或相关方略。编年条目可以为行动建立相对次序，却通常不保存劳动过程、物价变动、建筑材料或非精英的叙述。账本提示的研究线索为多语边疆档案、盟旗/寺院/绿洲地方史与帝国治理研究。事件所处条件为1750年冲突暴露既有权力安排的脆弱性。后续变化应限于制度条文与西藏地方实际政治运作需要并读藏汉材料。来源限制：以乾隆十六年（1751年）的同期文书为定位起点；官修记录、奏报或外交文本服务特定机构立场，具体日期、规模、地方执行与对手视角须以独立材料复核。
\hypertarget{civic:EQ-1752-TIBETAN-ADMINISTRATION}{}\paragraph{乾隆十七年（1752年）}驻藏系统与青海、川藏交通的军政协调继续调整。本事件的制度含义须同工程、运输、劳役、生态条件或知识传播的实际媒介结合，不能把一道制度或一部汇编当成地方实践本身。 核心取证为《清高宗实录》乾隆十七年相关条；军机处档、理藩院档或相关方略。编年条目可以为行动建立相对次序，却通常不保存劳动过程、物价变动、建筑材料或非精英的叙述。账本提示的研究线索为多语边疆档案、盟旗/寺院/绿洲地方史与帝国治理研究。事件所处条件为高原补给和多方政治关系限制中央命令的执行。后续变化应限于边务档案可显示决策链，不应替代地方社会证词。来源限制：以乾隆十七年（1752年）的同期文书为定位起点；官修记录、奏报或外交文本服务特定机构立场，具体日期、规模、地方执行与对手视角须以独立材料复核。
\hypertarget{civic:EQ-1753-SALT-ADMINISTRATION}{}\paragraph{乾隆十八年（1753年）}盐政、引岸和课额问题在中央与产销地区持续受到整顿。本事件的制度含义须同工程、运输、劳役、生态条件或知识传播的实际媒介结合，不能把一道制度或一部汇编当成地方实践本身。 核心取证为《清高宗实录》乾隆十八年相关条；户部奏销、盐政、河工或海关档案。行政文书显示的是机构如何把道路、港口、人口、宗教或案件转换为可处理的类别；没有后续县档或物质证据，不能由分类推定实际生活。账本提示的研究线索为财政账册、市场网络、灾害环境与区域经济研究。事件所处条件为盐课是重要财政来源，既有商人网络与官府规制交错。后续变化应限于制度命令不能直接推知盐价、私盐或商人收益。来源限制：以乾隆十八年（1753年）的同期文书为定位起点；官修记录、奏报或外交文本服务特定机构立场，具体日期、规模、地方执行与对手视角须以独立材料复核。
\hypertarget{civic:EQ-1754-AMURSANA-DEFECTION}{}\paragraph{乾隆十九年（1754年）}阿睦尔撒纳等准噶尔贵族向清廷归附，改变西北战争的联盟格局。本事件的制度含义须同工程、运输、劳役、生态条件或知识传播的实际媒介结合，不能把一道制度或一部汇编当成地方实践本身。 核心取证为《清高宗实录》乾隆十九年相关条；军机处档、理藩院档或相关方略。编年条目可以为行动建立相对次序，却通常不保存劳动过程、物价变动、建筑材料或非精英的叙述。账本提示的研究线索为多语边疆档案、盟旗/寺院/绿洲地方史与帝国治理研究。事件所处条件为准噶尔内部权力竞争加深。后续变化应限于清廷将其视为进军伊犁的机会，后续关系迅速破裂。来源限制：以乾隆十九年（1754年）的同期文书为定位起点；官修记录、奏报或外交文本服务特定机构立场，具体日期、规模、地方执行与对手视角须以独立材料复核。
\hypertarget{civic:EQ-1755-ILI-EXPEDITION}{}\paragraph{乾隆二十年（1755年）}清军分路进入伊犁，击溃达瓦齐政权并占领准噶尔核心地区。本事件的制度含义须同工程、运输、劳役、生态条件或知识传播的实际媒介结合，不能把一道制度或一部汇编当成地方实践本身。 核心取证为《清高宗实录》乾隆二十年相关条；军机处档、战事方略及地方奏报。编年条目可以为行动建立相对次序，却通常不保存劳动过程、物价变动、建筑材料或非精英的叙述。账本提示的研究线索为战争动员、军需、战场暴力与地方社会研究。事件所处条件为准噶尔内部裂解和清军长期准备相互配合。后续变化应限于占领并未立即稳定，阿睦尔撒纳再起使战争延续。来源限制：以乾隆二十年（1755年）的同期文书为定位起点；官修记录、奏报或外交文本服务特定机构立场，具体日期、规模、地方执行与对手视角须以独立材料复核。
\hypertarget{civic:EQ-1755-AMURSANA-REVOLT}{}\paragraph{乾隆二十年（1755年）}阿睦尔撒纳反清，准噶尔战后秩序重新陷入军事冲突。本事件的制度含义须同工程、运输、劳役、生态条件或知识传播的实际媒介结合，不能把一道制度或一部汇编当成地方实践本身。 核心取证为《清高宗实录》乾隆二十年相关条；军机处档、战事方略及地方奏报。编年条目可以为行动建立相对次序，却通常不保存劳动过程、物价变动、建筑材料或非精英的叙述。账本提示的研究线索为战争动员、军需、战场暴力与地方社会研究。事件所处条件为清廷对新归附力量的安置与阿睦尔撒纳的权力期待冲突。后续变化应限于清军扩大报复性军事行动，人口损失叙述须谨慎核验。来源限制：以乾隆二十年（1755年）的同期文书为定位起点；官修记录、奏报或外交文本服务特定机构立场，具体日期、规模、地方执行与对手视角须以独立材料复核。
\hypertarget{civic:EQ-1756-BANNER-CIVILIAN-REGISTRATION}{}\paragraph{乾隆二十一年（1756年）}清廷命令部分八旗副户口改归民籍，继续缓解旗籍供养压力。本事件的制度含义须同工程、运输、劳役、生态条件或知识传播的实际媒介结合，不能把一道制度或一部汇编当成地方实践本身。 核心取证为《清高宗实录》乾隆二十一年相关条；地方志、督抚奏折及地方档案。编年条目可以为行动建立相对次序，却通常不保存劳动过程、物价变动、建筑材料或非精英的叙述。账本提示的研究线索为地方档案、迁徙、宗教、族群与社会动员研究。事件所处条件为旗内人口与财政资源不匹配。后续变化应限于身份转变在地方具有不同的程序和实际结果。来源限制：以乾隆二十一年（1756年）的同期文书为定位起点；官修记录、奏报或外交文本服务特定机构立场，具体日期、规模、地方执行与对手视角须以独立材料复核。
\hypertarget{civic:EQ-1757-CANTON-TRADE-RESTRICTION}{}\paragraph{乾隆二十二年（1757年）}清廷将西洋贸易集中于广州口岸，形成后来所谓“一口通商”的制度框架。本事件的制度含义须同工程、运输、劳役、生态条件或知识传播的实际媒介结合，不能把一道制度或一部汇编当成地方实践本身。 核心取证为《清高宗实录》乾隆二十二年相关条；《筹办夷务始末》相关年卷与中外照会、条约文本。行政文书显示的是机构如何把道路、港口、人口、宗教或案件转换为可处理的类别；没有后续县档或物质证据，不能由分类推定实际生活。账本提示的研究线索为条约文本、外交档案、口岸社会与国际法实践研究。事件所处条件为沿海治安、关税和对外接触的管理需求上升。后续变化应限于广州贸易网络扩大，但禁限条文不等于所有海上往来消失。来源限制：以乾隆二十二年（1757年）的同期文书为定位起点；官修记录、奏报或外交文本服务特定机构立场，具体日期、规模、地方执行与对手视角须以独立材料复核。
\hypertarget{civic:EQ-1757-AMURSANA-DEATH}{}\paragraph{乾隆二十二年（1757年）}阿睦尔撒纳逃入俄境后死亡，清准战争的关键对手退出。本事件的制度含义须同工程、运输、劳役、生态条件或知识传播的实际媒介结合，不能把一道制度或一部汇编当成地方实践本身。 核心取证为《清高宗实录》乾隆二十二年相关条；军机处档、理藩院档或相关方略。编年条目可以为行动建立相对次序，却通常不保存劳动过程、物价变动、建筑材料或非精英的叙述。账本提示的研究线索为多语边疆档案、盟旗/寺院/绿洲地方史与帝国治理研究。事件所处条件为清军追击与俄清边境交涉交织。后续变化应限于其死亡与边境信息须结合俄方材料核对。来源限制：以乾隆二十二年（1757年）的同期文书为定位起点；官修记录、奏报或外交文本服务特定机构立场，具体日期、规模、地方执行与对手视角须以独立材料复核。
\hypertarget{civic:EQ-1758-KHOJA-REBELLION}{}\paragraph{乾隆二十三年（1758年）}大小和卓反清，清军由北路向天山南路推进。本事件的制度含义须同工程、运输、劳役、生态条件或知识传播的实际媒介结合，不能把一道制度或一部汇编当成地方实践本身。 核心取证为《清高宗实录》乾隆二十三年相关条；军机处档、战事方略及地方奏报。编年条目可以为行动建立相对次序，却通常不保存劳动过程、物价变动、建筑材料或非精英的叙述。账本提示的研究线索为战争动员、军需、战场暴力与地方社会研究。事件所处条件为伊犁占领后绿洲地区的权力安排尚未稳定。后续变化应限于清廷的军事胜利与绿洲社会的重新组织需分别追踪。来源限制：以乾隆二十三年（1758年）的同期文书为定位起点；官修记录、奏报或外交文本服务特定机构立场，具体日期、规模、地方执行与对手视角须以独立材料复核。
\hypertarget{civic:EQ-1759-KASHGAR-YARKAND-CAMPAIGN}{}\paragraph{乾隆二十四年（1759年）}清军攻入喀什噶尔、叶尔羌等地，和卓政权覆亡。本事件的制度含义须同工程、运输、劳役、生态条件或知识传播的实际媒介结合，不能把一道制度或一部汇编当成地方实践本身。 核心取证为《清高宗实录》乾隆二十四年相关条；军机处档、战事方略及地方奏报。编年条目可以为行动建立相对次序，却通常不保存劳动过程、物价变动、建筑材料或非精英的叙述。账本提示的研究线索为战争动员、军需、战场暴力与地方社会研究。事件所处条件为清军跨越天山并依赖绿洲补给和地方协力。后续变化应限于天山南路纳入清帝国的行政过程远晚于战役结束。来源限制：以乾隆二十四年（1759年）的同期文书为定位起点；官修记录、奏报或外交文本服务特定机构立场，具体日期、规模、地方执行与对手视角须以独立材料复核。
\subsection{1760—1769：区域网络、制度媒介与环境条件}
这些事件的先后可以由账本定位，但每一项影响都须回到具体区域、材料与行动者，不能被压缩为抽象的国家趋势。
\hypertarget{civic:EQ-1760-ILI-GENERAL}{}\paragraph{乾隆二十五年（1760年）}清廷设伊犁将军并展开驻防、屯田与移民安置。本事件的制度含义须同工程、运输、劳役、生态条件或知识传播的实际媒介结合，不能把一道制度或一部汇编当成地方实践本身。 核心取证为《清高宗实录》乾隆二十五年相关条；军机处档、理藩院档或相关方略。编年条目可以为行动建立相对次序，却通常不保存劳动过程、物价变动、建筑材料或非精英的叙述。账本提示的研究线索为多语边疆档案、盟旗/寺院/绿洲地方史与帝国治理研究。事件所处条件为西北战后需要长期军政管理和粮食供给。后续变化应限于驻防城市、屯田和地方社会间的关系不可简化为单向控制。来源限制：以乾隆二十五年（1760年）的同期文书为定位起点；官修记录、奏报或外交文本服务特定机构立场，具体日期、规模、地方执行与对手视角须以独立材料复核。
\hypertarget{civic:EQ-1761-XINJIANG-GARRISONS}{}\paragraph{乾隆二十六年（1761年）}清廷在新疆配置满汉蒙军队和绿营驻防，建立军政节点。本事件的制度含义须同工程、运输、劳役、生态条件或知识传播的实际媒介结合，不能把一道制度或一部汇编当成地方实践本身。 核心取证为《清高宗实录》乾隆二十六年相关条；军机处档、理藩院档或相关方略。编年条目可以为行动建立相对次序，却通常不保存劳动过程、物价变动、建筑材料或非精英的叙述。账本提示的研究线索为多语边疆档案、盟旗/寺院/绿洲地方史与帝国治理研究。事件所处条件为新征服地区交通遥远且人口与资源分布不均。后续变化应限于兵额与驻军实际存在需以档案、城址和账册互证。来源限制：以乾隆二十六年（1761年）的同期文书为定位起点；官修记录、奏报或外交文本服务特定机构立场，具体日期、规模、地方执行与对手视角须以独立材料复核。
\hypertarget{civic:EQ-1762-URUMQI-FOUNDATION}{}\paragraph{乾隆二十七年（1762年）}迪化等北疆城市的军政建设加快，成为屯田和交通枢纽。本事件的制度含义须同工程、运输、劳役、生态条件或知识传播的实际媒介结合，不能把一道制度或一部汇编当成地方实践本身。 核心取证为《清高宗实录》乾隆二十七年相关条；军机处档、理藩院档或相关方略。编年条目可以为行动建立相对次序，却通常不保存劳动过程、物价变动、建筑材料或非精英的叙述。账本提示的研究线索为多语边疆档案、盟旗/寺院/绿洲地方史与帝国治理研究。事件所处条件为伊犁—东疆之间需要稳定的补给与行政据点。后续变化应限于城市建设改变土地与水利使用，地方社会后果需另证。来源限制：以乾隆二十七年（1762年）的同期文书为定位起点；官修记录、奏报或外交文本服务特定机构立场，具体日期、规模、地方执行与对手视角须以独立材料复核。
\hypertarget{civic:EQ-1763-XINJIANG-SETTLEMENT}{}\paragraph{乾隆二十八年（1763年）}清廷继续在新疆安置屯田人口、军户和商业网络。本事件的制度含义须同工程、运输、劳役、生态条件或知识传播的实际媒介结合，不能把一道制度或一部汇编当成地方实践本身。 核心取证为《清高宗实录》乾隆二十八年相关条；地方志、督抚奏折及地方档案。编年条目可以为行动建立相对次序，却通常不保存劳动过程、物价变动、建筑材料或非精英的叙述。账本提示的研究线索为地方档案、迁徙、宗教、族群与社会动员研究。事件所处条件为驻防粮食、劳力和市场供给依赖移民与地方生产。后续变化应限于户籍与屯田册只能说明登记，不自动等于稳定定居。来源限制：以乾隆二十八年（1763年）的同期文书为定位起点；官修记录、奏报或外交文本服务特定机构立场，具体日期、规模、地方执行与对手视角须以独立材料复核。
\hypertarget{civic:EQ-1764-MANCHU-BANNER-RESETTLEMENT}{}\paragraph{乾隆二十九年（1764年）}部分八旗兵丁被调往西北驻防，旗籍家庭的迁移与供给成为边务问题。本事件的制度含义须同工程、运输、劳役、生态条件或知识传播的实际媒介结合，不能把一道制度或一部汇编当成地方实践本身。 核心取证为《清高宗实录》乾隆二十九年相关条；地方志、督抚奏折及地方档案。编年条目可以为行动建立相对次序，却通常不保存劳动过程、物价变动、建筑材料或非精英的叙述。账本提示的研究线索为地方档案、迁徙、宗教、族群与社会动员研究。事件所处条件为伊犁和北疆需要长期守备力量。后续变化应限于迁徙改变旗人家庭生活，不能只以军政文书描述。来源限制：以乾隆二十九年（1764年）的同期文书为定位起点；官修记录、奏报或外交文本服务特定机构立场，具体日期、规模、地方执行与对手视角须以独立材料复核。
\hypertarget{civic:EQ-1765-USH-REBELLION}{}\paragraph{乾隆三十年（1765年）}乌什爆发反清事件，清军镇压并重组当地治理。本事件的制度含义须同工程、运输、劳役、生态条件或知识传播的实际媒介结合，不能把一道制度或一部汇编当成地方实践本身。 核心取证为《清高宗实录》乾隆三十年相关条；军机处档、战事方略及地方奏报。编年条目可以为行动建立相对次序，却通常不保存劳动过程、物价变动、建筑材料或非精英的叙述。账本提示的研究线索为战争动员、军需、战场暴力与地方社会研究。事件所处条件为地方征发、驻军关系和绿洲权力矛盾累积。后续变化应限于官方定性和伤亡数字须与维吾尔文及地方材料交叉检验。来源限制：以乾隆三十年（1765年）的同期文书为定位起点；官修记录、奏报或外交文本服务特定机构立场，具体日期、规模、地方执行与对手视角须以独立材料复核。
\hypertarget{civic:EQ-1766-BURMA-WAR-BEGIN}{}\paragraph{乾隆三十一年（1766年）}清军对缅甸发动远征，西南边境战争展开。本事件的制度含义须同工程、运输、劳役、生态条件或知识传播的实际媒介结合，不能把一道制度或一部汇编当成地方实践本身。 核心取证为《清高宗实录》乾隆三十一年相关条；军机处档、战事方略及地方奏报。编年条目可以为行动建立相对次序，却通常不保存劳动过程、物价变动、建筑材料或非精英的叙述。账本提示的研究线索为战争动员、军需、战场暴力与地方社会研究。事件所处条件为边境贸易、土司关系和清缅政治冲突相互交织。后续变化应限于山地疾病、补给和气候削弱远征军的战斗力。来源限制：以乾隆三十一年（1766年）的同期文书为定位起点；官修记录、奏报或外交文本服务特定机构立场，具体日期、规模、地方执行与对手视角须以独立材料复核。
\hypertarget{civic:EQ-1767-BURMA-CAMPAIGN}{}\paragraph{乾隆三十二年（1767年）}清缅战争持续，清军在边境和山地路线中遭遇重大困难。本事件的制度含义须同工程、运输、劳役、生态条件或知识传播的实际媒介结合，不能把一道制度或一部汇编当成地方实践本身。 核心取证为《清高宗实录》乾隆三十二年相关条；军机处档、战事方略及地方奏报。编年条目可以为行动建立相对次序，却通常不保存劳动过程、物价变动、建筑材料或非精英的叙述。账本提示的研究线索为战争动员、军需、战场暴力与地方社会研究。事件所处条件为跨区域指挥与热带环境限制北方军队行动。后续变化应限于军事文书的战果叙述须与缅方材料并读。来源限制：以乾隆三十二年（1767年）的同期文书为定位起点；官修记录、奏报或外交文本服务特定机构立场，具体日期、规模、地方执行与对手视角须以独立材料复核。
\hypertarget{civic:EQ-1768-LITERARY-INQUISITION}{}\paragraph{乾隆三十三年（1768年）}乾隆朝以文字狱案件强化对著述、地方士人和政治语言的控制。本事件的制度含义须同工程、运输、劳役、生态条件或知识传播的实际媒介结合，不能把一道制度或一部汇编当成地方实践本身。 核心取证为《清高宗实录》乾隆三十三年相关条；内阁档、文集或报刊相关记录。编年条目可以为行动建立相对次序，却通常不保存劳动过程、物价变动、建筑材料或非精英的叙述。账本提示的研究线索为知识生产、教育、宗教与公共传播研究。事件所处条件为朝廷担忧地方文人以历史、族群或政论表达异议。后续变化应限于案件卷宗显示国家分类，不代表所有读书人的意图。来源限制：以乾隆三十三年（1768年）的同期文书为定位起点；官修记录、奏报或外交文本服务特定机构立场，具体日期、规模、地方执行与对手视角须以独立材料复核。
\hypertarget{civic:EQ-1768-SECOND-JINCHUAN-WAR}{}\paragraph{乾隆三十三年（1768年）}第二次金川战争开始，清廷投入更大兵力和资源。本事件的制度含义须同工程、运输、劳役、生态条件或知识传播的实际媒介结合，不能把一道制度或一部汇编当成地方实践本身。 核心取证为《清高宗实录》乾隆三十三年相关条；军机处档、战事方略及地方奏报。编年条目可以为行动建立相对次序，却通常不保存劳动过程、物价变动、建筑材料或非精英的叙述。账本提示的研究线索为战争动员、军需、战场暴力与地方社会研究。事件所处条件为第一次和议未能解决土司与清廷的权力冲突。后续变化应限于长期山地战带来巨额军费与地方劳役负担。来源限制：以乾隆三十三年（1768年）的同期文书为定位起点；官修记录、奏报或外交文本服务特定机构立场，具体日期、规模、地方执行与对手视角须以独立材料复核。
\hypertarget{civic:EQ-1769-BURMA-SETTLEMENT}{}\paragraph{乾隆三十四年（1769年）}清缅战争在反复交战后趋向停战和使节往来安排。本事件的制度含义须同工程、运输、劳役、生态条件或知识传播的实际媒介结合，不能把一道制度或一部汇编当成地方实践本身。 核心取证为《清高宗实录》乾隆三十四年相关条；《筹办夷务始末》相关年卷与中外照会、条约文本。行政文书显示的是机构如何把道路、港口、人口、宗教或案件转换为可处理的类别；没有后续县档或物质证据，不能由分类推定实际生活。账本提示的研究线索为条约文本、外交档案、口岸社会与国际法实践研究。事件所处条件为双方均承受远征成本与边境治理压力。后续变化应限于朝贡或使节语言不能直接等同单一的宗主权关系。来源限制：以乾隆三十四年（1769年）的同期文书为定位起点；官修记录、奏报或外交文本服务特定机构立场，具体日期、规模、地方执行与对手视角须以独立材料复核。
\hypertarget{civic:EQ-1769-JINCHUAN-ESCALATION}{}\paragraph{乾隆三十四年（1769年）}金川战场的道路、炮兵、粮运与将领调度不断升级。本事件的制度含义须同工程、运输、劳役、生态条件或知识传播的实际媒介结合，不能把一道制度或一部汇编当成地方实践本身。 核心取证为《清高宗实录》乾隆三十四年相关条；军机处档、战事方略及地方奏报。编年条目可以为行动建立相对次序，却通常不保存劳动过程、物价变动、建筑材料或非精英的叙述。账本提示的研究线索为战争动员、军需、战场暴力与地方社会研究。事件所处条件为山地堡寨体系使常规野战难以奏效。后续变化应限于清廷战后更重视边地交通与驻防，但其效果需地方材料检验。来源限制：以乾隆三十四年（1769年）的同期文书为定位起点；官修记录、奏报或外交文本服务特定机构立场，具体日期、规模、地方执行与对手视角须以独立材料复核。
\subsection{1770—1779：区域网络、制度媒介与环境条件}
这些事件的先后可以由账本定位，但每一项影响都须回到具体区域、材料与行动者，不能被压缩为抽象的国家趋势。
\hypertarget{civic:EQ-1770-SICHUAN-SUPPLY-BURDEN}{}\paragraph{乾隆三十五年（1770年）}金川军需加重四川及邻省的转运、民夫和财政负担。本事件的制度含义须同工程、运输、劳役、生态条件或知识传播的实际媒介结合，不能把一道制度或一部汇编当成地方实践本身。 核心取证为《清高宗实录》乾隆三十五年相关条；户部奏销、盐政、河工或海关档案。行政文书显示的是机构如何把道路、港口、人口、宗教或案件转换为可处理的类别；没有后续县档或物质证据，不能由分类推定实际生活。账本提示的研究线索为财政账册、市场网络、灾害环境与区域经济研究。事件所处条件为大规模驻军消耗粮食、牲畜和道路资源。后续变化应限于战时征发对村社和市场的影响不应被军功叙事遮蔽。来源限制：以乾隆三十五年（1770年）的同期文书为定位起点；官修记录、奏报或外交文本服务特定机构立场，具体日期、规模、地方执行与对手视角须以独立材料复核。
\hypertarget{civic:EQ-1771-TORGHUT-RETURN}{}\paragraph{乾隆三十六年（1771年）}土尔扈特部自伏尔加回归，清廷在伊犁等地安置部众。本事件的制度含义须同工程、运输、劳役、生态条件或知识传播的实际媒介结合，不能把一道制度或一部汇编当成地方实践本身。 核心取证为《清高宗实录》乾隆三十六年相关条；军机处档、理藩院档或相关方略。编年条目可以为行动建立相对次序，却通常不保存劳动过程、物价变动、建筑材料或非精英的叙述。账本提示的研究线索为多语边疆档案、盟旗/寺院/绿洲地方史与帝国治理研究。事件所处条件为俄国扩张、草原压力与部众迁徙决策共同作用。后续变化应限于安置文书不能掩盖迁徙损失和内部差异。来源限制：以乾隆三十六年（1771年）的同期文书为定位起点；官修记录、奏报或外交文本服务特定机构立场，具体日期、规模、地方执行与对手视角须以独立材料复核。
\hypertarget{civic:EQ-1772-SIKU-QUANSHU-ORDER}{}\paragraph{乾隆三十七年（1772年）}乾隆帝下令征集、编纂《四库全书》，建立全国性的图书征集机制。本事件的制度含义须同工程、运输、劳役、生态条件或知识传播的实际媒介结合，不能把一道制度或一部汇编当成地方实践本身。 核心取证为《清高宗实录》乾隆三十七年相关条；内阁档、文集或报刊相关记录。编年条目可以为行动建立相对次序，却通常不保存劳动过程、物价变动、建筑材料或非精英的叙述。账本提示的研究线索为知识生产、教育、宗教与公共传播研究。事件所处条件为朝廷希望整理典籍并审查政治敏感著作。后续变化应限于征书和禁毁并行，书目不能等于实际保存或流通。来源限制：以乾隆三十七年（1772年）的同期文书为定位起点；官修记录、奏报或外交文本服务特定机构立场，具体日期、规模、地方执行与对手视角须以独立材料复核。
\hypertarget{civic:EQ-1773-SICHUAN-RECRUITMENT}{}\paragraph{乾隆三十八年（1773年）}金川战事促使清廷继续在四川及邻省征调兵丁和夫役。本事件的制度含义须同工程、运输、劳役、生态条件或知识传播的实际媒介结合，不能把一道制度或一部汇编当成地方实践本身。 核心取证为《清高宗实录》乾隆三十八年相关条；军机处档、战事方略及地方奏报。编年条目可以为行动建立相对次序，却通常不保存劳动过程、物价变动、建筑材料或非精英的叙述。账本提示的研究线索为战争动员、军需、战场暴力与地方社会研究。事件所处条件为战场消耗超过初期预期。后续变化应限于兵役与夫役的地方分配存在显著不均。来源限制：以乾隆三十八年（1773年）的同期文书为定位起点；官修记录、奏报或外交文本服务特定机构立场，具体日期、规模、地方执行与对手视角须以独立材料复核。
\hypertarget{civic:EQ-1774-WANG-LUN-UPRISING}{}\paragraph{乾隆三十九年（1774年）}王伦在山东起事，清廷迅速调兵镇压。本事件的制度含义须同工程、运输、劳役、生态条件或知识传播的实际媒介结合，不能把一道制度或一部汇编当成地方实践本身。 核心取证为《清高宗实录》乾隆三十九年相关条；地方志、督抚奏折及地方档案。编年条目可以为行动建立相对次序，却通常不保存劳动过程、物价变动、建筑材料或非精英的叙述。账本提示的研究线索为地方档案、迁徙、宗教、族群与社会动员研究。事件所处条件为宗教结社、地方贫困和治安压力形成起事背景。后续变化应限于官府对“教门”的分类不应取代参与者多样诉求。来源限制：以乾隆三十九年（1774年）的同期文书为定位起点；官修记录、奏报或外交文本服务特定机构立场，具体日期、规模、地方执行与对手视角须以独立材料复核。
\hypertarget{civic:EQ-1775-SIKU-COLLECTION}{}\paragraph{乾隆四十年（1775年）}《四库全书》征书与审查在各省展开，地方藏书和文人网络受到影响。本事件的制度含义须同工程、运输、劳役、生态条件或知识传播的实际媒介结合，不能把一道制度或一部汇编当成地方实践本身。 核心取证为《清高宗实录》乾隆四十年相关条；内阁档、文集或报刊相关记录。编年条目可以为行动建立相对次序，却通常不保存劳动过程、物价变动、建筑材料或非精英的叙述。账本提示的研究线索为知识生产、教育、宗教与公共传播研究。事件所处条件为中央需要将地方文献纳入可控的知识秩序。后续变化应限于禁毁目录和地方执行范围应与存书、抄本和书信互证。来源限制：以乾隆四十年（1775年）的同期文书为定位起点；官修记录、奏报或外交文本服务特定机构立场，具体日期、规模、地方执行与对手视角须以独立材料复核。
\hypertarget{civic:EQ-1776-JINCHUAN-CONQUEST}{}\paragraph{乾隆四十一年（1776年）}清军攻克金川，第二次金川战争结束。本事件的制度含义须同工程、运输、劳役、生态条件或知识传播的实际媒介结合，不能把一道制度或一部汇编当成地方实践本身。 核心取证为《清高宗实录》乾隆四十一年相关条；军机处档、战事方略及地方奏报。编年条目可以为行动建立相对次序，却通常不保存劳动过程、物价变动、建筑材料或非精英的叙述。账本提示的研究线索为战争动员、军需、战场暴力与地方社会研究。事件所处条件为长期工程、火器和大规模转运最终压倒山地防御。后续变化应限于战后改土归流与人口迁移的过程不应由“平定”一词遮蔽。来源限制：以乾隆四十一年（1776年）的同期文书为定位起点；官修记录、奏报或外交文本服务特定机构立场，具体日期、规模、地方执行与对手视角须以独立材料复核。
\hypertarget{civic:EQ-1777-JINCHUAN-ADMINISTRATIVE-RESET}{}\paragraph{乾隆四十二年（1777年）}清廷在金川地区推行改土归流、设治和驻防安排。本事件的制度含义须同工程、运输、劳役、生态条件或知识传播的实际媒介结合，不能把一道制度或一部汇编当成地方实践本身。 核心取证为《清高宗实录》乾隆四十二年相关条；军机处档、理藩院档或相关方略。编年条目可以为行动建立相对次序，却通常不保存劳动过程、物价变动、建筑材料或非精英的叙述。账本提示的研究线索为多语边疆档案、盟旗/寺院/绿洲地方史与帝国治理研究。事件所处条件为军事胜利后需要将战果转为持续行政控制。后续变化应限于地方秩序重组伴随土地、劳役和宗教网络变化。来源限制：以乾隆四十二年（1777年）的同期文书为定位起点；官修记录、奏报或外交文本服务特定机构立场，具体日期、规模、地方执行与对手视角须以独立材料复核。
\hypertarget{civic:EQ-1778-SICHUAN-TAX-RECOVERY}{}\paragraph{乾隆四十三年（1778年）}金川战后四川财政和民力恢复成为地方官奏报的重要议题。本事件的制度含义须同工程、运输、劳役、生态条件或知识传播的实际媒介结合，不能把一道制度或一部汇编当成地方实践本身。 核心取证为《清高宗实录》乾隆四十三年相关条；户部奏销、盐政、河工或海关档案。行政文书显示的是机构如何把道路、港口、人口、宗教或案件转换为可处理的类别；没有后续县档或物质证据，不能由分类推定实际生活。账本提示的研究线索为财政账册、市场网络、灾害环境与区域经济研究。事件所处条件为长期战争造成钱粮、道路和人口损耗。后续变化应限于官方恢复叙事须同地方账册和市场材料核对。来源限制：以乾隆四十三年（1778年）的同期文书为定位起点；官修记录、奏报或外交文本服务特定机构立场，具体日期、规模、地方执行与对手视角须以独立材料复核。
\hypertarget{civic:EQ-1779-VIETNAM-INTERVENTION}{}\paragraph{乾隆四十四年（1779年）}清廷介入安南政局并派兵南下，孙士毅军一度进入升龙。本事件的制度含义须同工程、运输、劳役、生态条件或知识传播的实际媒介结合，不能把一道制度或一部汇编当成地方实践本身。 核心取证为《清高宗实录》乾隆四十四年相关条；军机处档、战事方略及地方奏报。编年条目可以为行动建立相对次序，却通常不保存劳动过程、物价变动、建筑材料或非精英的叙述。账本提示的研究线索为战争动员、军需、战场暴力与地方社会研究。事件所处条件为黎朝复辟请求与区域权力竞争相交。后续变化应限于清军行动受当地政治与补给限制，战果叙述有待中越材料互校。来源限制：以乾隆四十四年（1779年）的同期文书为定位起点；官修记录、奏报或外交文本服务特定机构立场，具体日期、规模、地方执行与对手视角须以独立材料复核。
\subsection{1780—1789：区域网络、制度媒介与环境条件}
这些事件的先后可以由账本定位，但每一项影响都须回到具体区域、材料与行动者，不能被压缩为抽象的国家趋势。
\hypertarget{civic:EQ-1780-VIETNAM-WITHDRAWAL}{}\paragraph{乾隆四十五年（1780年）}清军撤离安南，双方通过使节和册封语言恢复关系。本事件的制度含义须同工程、运输、劳役、生态条件或知识传播的实际媒介结合，不能把一道制度或一部汇编当成地方实践本身。 核心取证为《清高宗实录》乾隆四十五年相关条；《筹办夷务始末》相关年卷与中外照会、条约文本。行政文书显示的是机构如何把道路、港口、人口、宗教或案件转换为可处理的类别；没有后续县档或物质证据，不能由分类推定实际生活。账本提示的研究线索为条约文本、外交档案、口岸社会与国际法实践研究。事件所处条件为军事行动未能建立稳定的复辟政权。后续变化应限于外交礼仪提供缓和机制，不等于消除区域竞争。来源限制：以乾隆四十五年（1780年）的同期文书为定位起点；官修记录、奏报或外交文本服务特定机构立场，具体日期、规模、地方执行与对手视角须以独立材料复核。
\hypertarget{civic:EQ-1781-GANSU-MUSLIM-UPRISING}{}\paragraph{乾隆四十六年（1781年）}甘肃爆发回民起事，清军实施镇压与后续治理。本事件的制度含义须同工程、运输、劳役、生态条件或知识传播的实际媒介结合，不能把一道制度或一部汇编当成地方实践本身。 核心取证为《清高宗实录》乾隆四十六年相关条；军机处档、战事方略及地方奏报。编年条目可以为行动建立相对次序，却通常不保存劳动过程、物价变动、建筑材料或非精英的叙述。账本提示的研究线索为战争动员、军需、战场暴力与地方社会研究。事件所处条件为地方纠纷、行政失当和族群分类被动员为冲突语言。后续变化应限于伤亡、动机与群体边界必须以地方、多语材料审慎处理。来源限制：以乾隆四十六年（1781年）的同期文书为定位起点；官修记录、奏报或外交文本服务特定机构立场，具体日期、规模、地方执行与对手视角须以独立材料复核。
\hypertarget{civic:EQ-1782-GANSU-POSTWAR-GOVERNANCE}{}\paragraph{乾隆四十七年（1782年）}清廷在甘肃起事后调整驻兵、司法和地方控制措施。本事件的制度含义须同工程、运输、劳役、生态条件或知识传播的实际媒介结合，不能把一道制度或一部汇编当成地方实践本身。 核心取证为《清高宗实录》乾隆四十七年相关条；地方志、督抚奏折及地方档案。编年条目可以为行动建立相对次序，却通常不保存劳动过程、物价变动、建筑材料或非精英的叙述。账本提示的研究线索为地方档案、迁徙、宗教、族群与社会动员研究。事件所处条件为战后需要恢复交通、赋役与社区关系。后续变化应限于官方治安叙事不能替代地方生计和迁徙记录。来源限制：以乾隆四十七年（1782年）的同期文书为定位起点；官修记录、奏报或外交文本服务特定机构立场，具体日期、规模、地方执行与对手视角须以独立材料复核。
\hypertarget{civic:EQ-1783-SALT-MONOPOLY-REVIEW}{}\paragraph{乾隆四十八年（1783年）}清廷继续讨论盐政弊端、商人负担和课额征解。本事件的制度含义须同工程、运输、劳役、生态条件或知识传播的实际媒介结合，不能把一道制度或一部汇编当成地方实践本身。 核心取证为《清高宗实录》乾隆四十八年相关条；户部奏销、盐政、河工或海关档案。行政文书显示的是机构如何把道路、港口、人口、宗教或案件转换为可处理的类别；没有后续县档或物质证据，不能由分类推定实际生活。账本提示的研究线索为财政账册、市场网络、灾害环境与区域经济研究。事件所处条件为盐课、运销网络与地方财政紧密相连。后续变化应限于改革方案须与实际盐价、私盐和账册区分。来源限制：以乾隆四十八年（1783年）的同期文书为定位起点；官修记录、奏报或外交文本服务特定机构立场，具体日期、规模、地方执行与对手视角须以独立材料复核。
\hypertarget{civic:EQ-1785-TAIWAN-GARRISON}{}\paragraph{乾隆五十年（1785年）}清廷调整台湾驻防、渡台与地方治安安排。本事件的制度含义须同工程、运输、劳役、生态条件或知识传播的实际媒介结合，不能把一道制度或一部汇编当成地方实践本身。 核心取证为《清高宗实录》乾隆五十年相关条；军机处档、理藩院档或相关方略。编年条目可以为行动建立相对次序，却通常不保存劳动过程、物价变动、建筑材料或非精英的叙述。账本提示的研究线索为多语边疆档案、盟旗/寺院/绿洲地方史与帝国治理研究。事件所处条件为海峡调兵和岛内治理成本持续上升。后续变化应限于制度设计与港口、村社的实际执行需要分开考察。来源限制：以乾隆五十年（1785年）的同期文书为定位起点；官修记录、奏报或外交文本服务特定机构立场，具体日期、规模、地方执行与对手视角须以独立材料复核。
\hypertarget{civic:EQ-1786-LIN-SHUANGWEN-UPRISING}{}\paragraph{乾隆五十一年（1786年）}林爽文起事在台湾爆发并迅速扩展。本事件的制度含义须同工程、运输、劳役、生态条件或知识传播的实际媒介结合，不能把一道制度或一部汇编当成地方实践本身。 核心取证为《清高宗实录》乾隆五十一年相关条；军机处档、战事方略及地方奏报。编年条目可以为行动建立相对次序，却通常不保存劳动过程、物价变动、建筑材料或非精英的叙述。账本提示的研究线索为战争动员、军需、战场暴力与地方社会研究。事件所处条件为地方权力竞争、赋役与治安危机叠加。后续变化应限于清廷须跨海调兵，地方参与者的目标并不一致。来源限制：以乾隆五十一年（1786年）的同期文书为定位起点；官修记录、奏报或外交文本服务特定机构立场，具体日期、规模、地方执行与对手视角须以独立材料复核。
\hypertarget{civic:EQ-1786-FIRST-GURKHA-WAR}{}\paragraph{乾隆五十一年（1786年）}廓尔喀军进入西藏边境，清廷开始处理高原军事危机。本事件的制度含义须同工程、运输、劳役、生态条件或知识传播的实际媒介结合，不能把一道制度或一部汇编当成地方实践本身。 核心取证为《清高宗实录》乾隆五十一年相关条；军机处档、理藩院档或相关方略。编年条目可以为行动建立相对次序，却通常不保存劳动过程、物价变动、建筑材料或非精英的叙述。账本提示的研究线索为多语边疆档案、盟旗/寺院/绿洲地方史与帝国治理研究。事件所处条件为贸易、货币与边界纠纷恶化。后续变化应限于清藏与廓尔喀材料对冲突起因有不同叙述。来源限制：以乾隆五十一年（1786年）的同期文书为定位起点；官修记录、奏报或外交文本服务特定机构立场，具体日期、规模、地方执行与对手视角须以独立材料复核。
\hypertarget{civic:EQ-1787-TAIWAN-CAMPAIGN}{}\paragraph{乾隆五十二年（1787年）}福康安等率军赴台镇压林爽文起事，海峡运输成为关键。本事件的制度含义须同工程、运输、劳役、生态条件或知识传播的实际媒介结合，不能把一道制度或一部汇编当成地方实践本身。 核心取证为《清高宗实录》乾隆五十二年相关条；军机处档、战事方略及地方奏报。编年条目可以为行动建立相对次序，却通常不保存劳动过程、物价变动、建筑材料或非精英的叙述。账本提示的研究线索为战争动员、军需、战场暴力与地方社会研究。事件所处条件为岛内守军不足且起事波及多地。后续变化应限于军事镇压与地方招抚、清算应分别记录。来源限制：以乾隆五十二年（1787年）的同期文书为定位起点；官修记录、奏报或外交文本服务特定机构立场，具体日期、规模、地方执行与对手视角须以独立材料复核。
\hypertarget{civic:EQ-1788-LIN-SHUANGWEN-DEFEAT}{}\paragraph{乾隆五十三年（1788年）}清军击败林爽文，台湾起事基本结束。本事件的制度含义须同工程、运输、劳役、生态条件或知识传播的实际媒介结合，不能把一道制度或一部汇编当成地方实践本身。 核心取证为《清高宗实录》乾隆五十三年相关条；军机处档、战事方略及地方奏报。编年条目可以为行动建立相对次序，却通常不保存劳动过程、物价变动、建筑材料或非精英的叙述。账本提示的研究线索为战争动员、军需、战场暴力与地方社会研究。事件所处条件为增援军队、地方协力和起事内部困难共同影响战局。后续变化应限于战后处置、伤亡和地方秩序重建不能由官书凯旋叙事概括。来源限制：以乾隆五十三年（1788年）的同期文书为定位起点；官修记录、奏报或外交文本服务特定机构立场，具体日期、规模、地方执行与对手视角须以独立材料复核。
\hypertarget{civic:EQ-1788-SECOND-GURKHA-WAR}{}\paragraph{乾隆五十三年（1788年）}廓尔喀再次入侵西藏，清军组织更大规模高原远征。本事件的制度含义须同工程、运输、劳役、生态条件或知识传播的实际媒介结合，不能把一道制度或一部汇编当成地方实践本身。 核心取证为《清高宗实录》乾隆五十三年相关条；军机处档、战事方略及地方奏报。编年条目可以为行动建立相对次序，却通常不保存劳动过程、物价变动、建筑材料或非精英的叙述。账本提示的研究线索为战争动员、军需、战场暴力与地方社会研究。事件所处条件为前次安排未化解边境贸易与政治矛盾。后续变化应限于军需和高原路线使战争超出拉萨单一政治中心。来源限制：以乾隆五十三年（1788年）的同期文书为定位起点；官修记录、奏报或外交文本服务特定机构立场，具体日期、规模、地方执行与对手视角须以独立材料复核。
\hypertarget{civic:EQ-1789-TIBET-NEPAL-CAMPAIGN}{}\paragraph{乾隆五十四年（1789年）}清军越过喜马拉雅方向追击廓尔喀，边境协商与军事行动并行。本事件的制度含义须同工程、运输、劳役、生态条件或知识传播的实际媒介结合，不能把一道制度或一部汇编当成地方实践本身。 核心取证为《清高宗实录》乾隆五十四年相关条；军机处档、战事方略及地方奏报。编年条目可以为行动建立相对次序，却通常不保存劳动过程、物价变动、建筑材料或非精英的叙述。账本提示的研究线索为战争动员、军需、战场暴力与地方社会研究。事件所处条件为清廷意在重建西藏边防威慑。后续变化应限于行军路线、战果和伤亡须以藏、尼泊尔及清方材料核验。来源限制：以乾隆五十四年（1789年）的同期文书为定位起点；官修记录、奏报或外交文本服务特定机构立场，具体日期、规模、地方执行与对手视角须以独立材料复核。
\subsection{1790—1799：区域网络、制度媒介与环境条件}
这些事件的先后可以由账本定位，但每一项影响都须回到具体区域、材料与行动者，不能被压缩为抽象的国家趋势。
\hypertarget{civic:EQ-1790-IMPERIAL-REGULATIONS-TIBET}{}\paragraph{乾隆五十五年（1790年）}清廷制定整顿西藏事务的制度性安排，强化驻藏大臣和金瓶掣签等机制。本事件的制度含义须同工程、运输、劳役、生态条件或知识传播的实际媒介结合，不能把一道制度或一部汇编当成地方实践本身。 核心取证为《清高宗实录》乾隆五十五年相关条；军机处档、理藩院档或相关方略。编年条目可以为行动建立相对次序，却通常不保存劳动过程、物价变动、建筑材料或非精英的叙述。账本提示的研究线索为多语边疆档案、盟旗/寺院/绿洲地方史与帝国治理研究。事件所处条件为廓尔喀战争暴露既有治理安排的漏洞。后续变化应限于制度条文与其后宗教、地方政治实践不能直接画等号。来源限制：以乾隆五十五年（1790年）的同期文书为定位起点；官修记录、奏报或外交文本服务特定机构立场，具体日期、规模、地方执行与对手视角须以独立材料复核。
\hypertarget{civic:EQ-1791-WHITE-LOTUS-PRECURSORS}{}\paragraph{乾隆五十六年（1791年）}川楚陕基层社会的教门网络和地方压力成为白莲教战争的背景。本事件的制度含义须同工程、运输、劳役、生态条件或知识传播的实际媒介结合，不能把一道制度或一部汇编当成地方实践本身。 核心取证为《清高宗实录》乾隆五十六年相关条；地方志、督抚奏折及地方档案。编年条目可以为行动建立相对次序，却通常不保存劳动过程、物价变动、建筑材料或非精英的叙述。账本提示的研究线索为地方档案、迁徙、宗教、族群与社会动员研究。事件所处条件为人口、土地、赋役和地方治安矛盾长期积累。后续变化应限于官府分类为“邪教”的材料须与社区网络分开分析。来源限制：以乾隆五十六年（1791年）的同期文书为定位起点；官修记录、奏报或外交文本服务特定机构立场，具体日期、规模、地方执行与对手视角须以独立材料复核。
\hypertarget{civic:EQ-1792-GURKHA-SETTLEMENT}{}\paragraph{乾隆五十七年（1792年）}清廷与廓尔喀战争结束后处理使节、边防和西藏军政善后。本事件的制度含义须同工程、运输、劳役、生态条件或知识传播的实际媒介结合，不能把一道制度或一部汇编当成地方实践本身。 核心取证为《清高宗实录》乾隆五十七年相关条；《筹办夷务始末》相关年卷与中外照会、条约文本。行政文书显示的是机构如何把道路、港口、人口、宗教或案件转换为可处理的类别；没有后续县档或物质证据，不能由分类推定实际生活。账本提示的研究线索为条约文本、外交档案、口岸社会与国际法实践研究。事件所处条件为远征成本与边境稳定需要可持续安排。后续变化应限于朝贡表述与边境实际互动并不完全一致。来源限制：以乾隆五十七年（1792年）的同期文书为定位起点；官修记录、奏报或外交文本服务特定机构立场，具体日期、规模、地方执行与对手视角须以独立材料复核。
\hypertarget{civic:EQ-1793-MACARTNEY-MISSION}{}\paragraph{乾隆五十八年（1793年）}马戛尔尼使团抵达清廷，围绕贸易、礼仪和外交关系交涉。本事件的制度含义须同工程、运输、劳役、生态条件或知识传播的实际媒介结合，不能把一道制度或一部汇编当成地方实践本身。 核心取证为《清高宗实录》乾隆五十八年相关条；《筹办夷务始末》相关年卷与中外照会、条约文本。行政文书显示的是机构如何把道路、港口、人口、宗教或案件转换为可处理的类别；没有后续县档或物质证据，不能由分类推定实际生活。账本提示的研究线索为条约文本、外交档案、口岸社会与国际法实践研究。事件所处条件为英国对华贸易与全球帝国扩张需求增长。后续变化应限于使团记录具有强烈英方视角，礼仪争议不应简化为“闭关”。来源限制：以乾隆五十八年（1793年）的同期文书为定位起点；官修记录、奏报或外交文本服务特定机构立场，具体日期、规模、地方执行与对手视角须以独立材料复核。
\hypertarget{civic:EQ-1794-WHITE-LOTUS-OUTBREAK}{}\paragraph{乾隆五十九年（1794年）}白莲教相关起事在川楚陕地区扩展，长期战争开始。本事件的制度含义须同工程、运输、劳役、生态条件或知识传播的实际媒介结合，不能把一道制度或一部汇编当成地方实践本身。 核心取证为《清高宗实录》乾隆五十九年相关条；军机处档、战事方略及地方奏报。编年条目可以为行动建立相对次序，却通常不保存劳动过程、物价变动、建筑材料或非精英的叙述。账本提示的研究线索为战争动员、军需、战场暴力与地方社会研究。事件所处条件为地方生计压力、宗教网络和治理失灵相互作用。后续变化应限于“教匪”标签不能替代对不同参与群体的分析。来源限制：以乾隆五十九年（1794年）的同期文书为定位起点；官修记录、奏报或外交文本服务特定机构立场，具体日期、规模、地方执行与对手视角须以独立材料复核。
\hypertarget{civic:EQ-1795-QIANLONG-ABDICATION}{}\paragraph{乾隆六十年（1795年）}乾隆帝禅位于嘉庆帝，自为太上皇，仍保留重要政治影响。本事件的制度含义须同工程、运输、劳役、生态条件或知识传播的实际媒介结合，不能把一道制度或一部汇编当成地方实践本身。 核心取证为《清高宗实录》乾隆六十年相关条；宫中朱批奏折、内阁大库题本与《清史稿》相关志传。行政文书显示的是机构如何把道路、港口、人口、宗教或案件转换为可处理的类别；没有后续县档或物质证据，不能由分类推定实际生活。账本提示的研究线索为宫廷政治、制度运作与中央—地方信息链研究。事件所处条件为乾隆希望履行不逾祖父在位年数的承诺。后续变化应限于嘉庆初年权力交接受太上皇与和珅网络制约。来源限制：以乾隆六十年（1795年）的同期文书为定位起点；官修记录、奏报或外交文本服务特定机构立场，具体日期、规模、地方执行与对手视角须以独立材料复核。
